\chapter*{Resumen}
La anemia falciforme continúa siendo una condición hemoglobinopática de alto impacto social y clínico, especialmente en regiones donde el acceso a laboratorios especializados resulta limitado. Las metodologías diagnósticas estándar cromatografía líquida de alta resolución (HPLC) o la electroforesis dependen de equipamiento costoso, procesos de tinción y personal altamente entrenado, lo que retrasa campañas de cribado masivo en comunidades vulnerables. En respuesta a esta brecha, el presente trabajo de grado propone un sistema de clasificación automatizado que analiza hologramas generados mediante microscopía digital sin lentes. La hipótesis sostiene que, aun sin reconstruir el campo óptico, es posible discriminar entre glóbulos rojos sanos y células falciformes si se combinan descriptores de textura, forma y frecuencia con modelos de aprendizaje automático clásicos y un módulo de detección de anomalías. Permitiendo así eliminar la parte más costosa computacionalmente del proceso de clasificación que es la reconstrucción en fase

El estudio parte de un conjunto de 304 hologramas rotulados (151 sanos y 153 enfermos) y de un repositorio de validación externa compuesto por hologramas de organismos ajenos a la sangre. Ambos datasets fueron proporcionados por el semillero de óptica y fotónica de la Universidad EAFIT. La metodología abarca la unificación automatizada de bases heterogéneas, la extracción de 84 características físicas y estadísticas, la selección de atributos relevantes y la construcción de un ensamble de clasificadores compuesto por regresión logística, máquinas de soporte vectorial, bosques aleatorios y gradiente por potenciación. Se integró, además, un detector de anomalías basado en la distancia de Mahalanobis para identificar muestras fuera de dominio y complementar la interpretación de las decisiones del modelo principal.

El pipeline se implementó en una interfaz de línea de comandos reproducible con un único perfil operativo que encapsula la unificación de datos, el entrenamiento del ensamble y la detección de anomalías. Las campañas experimentales alcanzaron una exactitud de prueba del 95,1\% y un área bajo la curva ROC de 0,989, asistidas por validación cruzada estratificada, análisis Leave-One-Out y estimaciones bootstrap de incertidumbre. Asimismo, la detección de anomalías etiquetó correctamente el 100\% de los hologramas no sanguíneos. Los resultados demuestran que la sinergia entre óptica computacional y aprendizaje automático interpretable habilita soluciones de diagnóstico portátiles, transparentes y escalables.

\textbf{Palabras clave:} anemia falciforme, microscopía holográfica sin lente, aprendizaje automático interpretable, detección de anomalías, análisis de patrones.
